\documentclass[11pt]{article}

\usepackage[margin=1in]{geometry}
\usepackage{amsmath,amssymb}
\usepackage{array}
\usepackage{booktabs}
\usepackage{caption}
\usepackage{hyperref}
\emergencystretch=2em

\title{Teacher-Anchored and Statistical Calibration for Drift-Adaptive Affine GKP Decoding}
\author{}
\date{}

\begin{document}
\maketitle

\begin{abstract}
Approximate Gottesman--Kitaev--Preskill (GKP) quantum error correction
encodes a qubit into an oscillator and uses continuous-valued modular
syndromes to infer displacement errors.  In realistic finite-energy GKP
states, syndrome statistics are shaped by state-preparation width,
measurement noise, circuit-level imperfections, and slow calibration drift.
This work studies a two-timescale decoding architecture for such settings.  A
low-latency fast loop applies an affine correction
\(\Delta_t=K_t s_t+b_t\), while a slower calibration loop estimates the current
effective noise state from recent syndrome histograms and updates the runtime
parameters.  This keeps the per-shot path compatible with fixed-point hardware
execution while allowing the decoder to adapt to nonstationary noise.

The architecture is designed around a narrow and deployment-oriented role for
learning.  A classical teacher estimator provides an interpretable baseline
for the affine parameters, and a lightweight CNN predicts only a bounded
residual calibration term, currently focused on the bias \(b_t\).  This differs
from end-to-end neural decoders and from outer-code soft-information decoders:
the learned module updates a low-dimensional physical-layer calibration
surface, whereas the real-time path remains a deterministic affine rule.
Software-HIL experiments over four drift scenarios show that adaptive affine
calibration improves over filtering baselines, that histogram-delta features
are important for the learned residual branch, and that a statistical
calibration variant substantially reduces the logical-error proxy across the
same drift suite.  A local sensitivity study further indicates that the
statistical-calibration advantage is not tied to a single narrow parameter
choice.
\end{abstract}

\tableofcontents

\section{Introduction}

Bosonic codes protect quantum information by encoding a logical qubit into the
Hilbert space of an oscillator.  The GKP code is a central example: ideal code
states form a periodic lattice in phase space, so small shift errors can be
diagnosed through modular syndrome measurements and corrected by displacement
operations \cite{gkp2001,grimsmo2021}.  Unlike binary stabilizer syndromes, GKP
syndromes are analog.  Their continuous values carry information not only
about the likely displacement, but also about how close the correction is to a
lattice decision boundary.

The practical challenge is that physical GKP states are approximate
finite-energy states.  Their comb peaks have finite width and an envelope, and
syndrome measurements are further affected by measurement inefficiency, noisy
auxiliary states, circuit-level imperfections, oscillator loss, and calibration
error \cite{grimsmo2021,roy2025}.  Consequently, the effective noise model seen
by the decoder can change over time.  A fixed syndrome-to-correction map may be
well tuned for one operating point but mismatched after bias, variance, or
covariance drift.

Recent GKP and bosonic-decoding work makes clear that analog soft information
is valuable.  Surface-GKP and concatenated bosonic-QLDPC decoders use
continuous syndrome information to set matching weights, belief-propagation
messages, or effective outer-code priors \cite{noh2020,noh2022,raveendran2022,
berent2024,borah2025}.  At the same time, modern QEC decoder systems show that
learned modules are most compelling when their latency role and integration
boundary are explicit \cite{chamberland2026,stein2026}.  Hardware and
calibration-oriented studies similarly emphasize that decoder priors must be
kept aligned with device statistics \cite{spitz2018,wagner2021,dgr2023,
chen2022,sivak2024}.

This work applies that systems principle at the physical GKP correction layer.
The fast loop is intentionally simple:
\begin{equation}
    \Delta_t = K_t s_t + b_t ,
\end{equation}
where \(s_t\) is the measured two-quadrature GKP syndrome and
\((K_t,b_t)\) are quantized runtime parameters.  The slow loop observes recent
syndrome histograms, estimates an effective noise state, and updates
\((K_t,b_t)\) at a lower rate.  The design therefore separates two concerns:
the per-shot correction path is a bounded affine operation, while the
nonstationary estimation problem is handled asynchronously through calibrated
parameter updates.

The main advantage of this formulation is its combination of physical
structure, adaptation, and deployment discipline.  Compared with a fixed
affine decoder, the parameters can track drift.  Compared with a full neural
decoder, the online computation is smaller and more predictable.  Compared
with outer-code soft-information methods, the adaptation target is a
physical-layer GKP displacement rule rather than a matching graph or LDPC
message schedule.  Compared with calibration-conditioned neural decoders, the
slow context does not enlarge the per-shot path; it only changes the affine
parameters staged into the runtime bank.

\section{Summary of Contributions}

The work can be organized around four contributions.

\paragraph{1. A two-timescale adaptive affine formulation for approximate GKP
decoding.}
The per-shot correction path is formulated as a bounded affine estimator
\(\Delta_t=K_t s_t+b_t\).  The expensive adaptation problem is moved into a
slow loop that estimates the current effective noise state from recent syndrome
statistics.  This links the continuous-syndrome GKP picture to a runtime
contract that can be implemented as a quantized matrix-vector operation.

\paragraph{2. A teacher-anchored residual calibration strategy.}
A classical teacher estimator produces a stable baseline
\((K_t^{\rm teacher},b_t^{\rm teacher})\).  The CNN predicts a small residual,
\begin{equation}
    \delta b_t=f_\phi(H_{t-c+1:t},\Delta H_t,z_t^{\rm teacher}),
\end{equation}
and the committed parameters are
\begin{equation}
    K_t=K_t^{\rm teacher}, \qquad
    b_t={\rm EMA}\!\left(b_t^{\rm teacher}+\delta b_t\right).
\end{equation}
Learning is therefore used to calibrate a low-dimensional control surface, not
to replace the GKP decoder end-to-end.

\paragraph{3. A deployment-aware runtime architecture.}
The architecture exposes fixed-point quantization, clipping, saturation
diagnostics, parameter smoothing, stale-parameter behavior, and double-buffered
stage-and-commit updates.  These constraints become measurable objects in the
benchmark rather than assumptions added after the algorithm is designed.

\paragraph{4. A benchmark suite for adaptive affine calibration.}
The numerical study compares filtering baselines, learned residual variants,
and statistical calibration variants under four drift scenarios.  The results
separate three questions: whether adaptive affine decoding improves over
filtering baselines, which learned residual features matter, and whether a
simple statistical calibration rule can provide a strong non-neural comparator
inside the same runtime contract.

\section{Brief Review of the GKP Code}

\subsection{Ideal and approximate code states}

The square-lattice GKP code encodes a qubit into an oscillator by using
periodic structure in the \(q\) and \(p\) quadratures \cite{gkp2001}.  In the
convention used here, the lattice constant is
\begin{equation}
    \lambda=\sqrt{2\pi}.
\end{equation}
An ideal logical state may be represented heuristically as an infinite comb,
for example
\begin{equation}
    |\bar 0\rangle \propto \sum_{n\in\mathbb{Z}} |n\lambda\rangle_q .
\end{equation}
Ideal comb states are not physical because they require infinite energy.
Approximate GKP states replace the infinitely sharp comb with finitely squeezed
peaks and an envelope.  Finite-energy structure is central to practical GKP
QEC rather than a minor implementation detail \cite{grimsmo2021}: it
determines the intrinsic syndrome uncertainty and contributes an effective
displacement-like noise.

\subsection{Syndrome measurement as modular displacement information}

Let the accumulated phase-space displacement error before correction be
\begin{equation}
    e_t=\begin{bmatrix}e_{q,t}\\e_{p,t}\end{bmatrix}.
\end{equation}
An ideal GKP syndrome observes this displacement modulo the lattice:
\begin{equation}
    s_t = e_t \bmod \lambda,
    \qquad
    s_{q,t},s_{p,t}\in[-\lambda/2,\lambda/2).
\end{equation}
In a finite-energy and noisy-measurement setting, the measured syndrome is
more accurately described as
\begin{equation}
    \tilde{s}_t={\rm mod}(e_t,\lambda)+\eta_t^{\rm meas}+\eta_t^{\rm GKP}.
\end{equation}
Here \(\eta_t^{\rm meas}\) represents measurement and auxiliary-state
contributions, while \(\eta_t^{\rm GKP}\) summarizes finite-squeezing
uncertainty.  The decoder therefore observes a noisy modular representative,
not the absolute displacement.

\subsection{Local affine decoding and its limits}

The modulo structure makes exact GKP decoding nonlinear and branch dependent.
Near a single lattice branch, however, a Gaussian approximation gives a useful
local model.  If \(e\) and \(s\) are jointly Gaussian,
\begin{equation}
\begin{bmatrix}e\\s\end{bmatrix}
\sim
\mathcal{N}
\left(
\begin{bmatrix}\mu_e\\\mu_s\end{bmatrix},
\begin{bmatrix}
\Sigma_{ee} & \Sigma_{es}\\
\Sigma_{se} & \Sigma_{ss}
\end{bmatrix}
\right),
\end{equation}
then the linear-MMSE estimate is
\begin{equation}
    \hat e=\mu_e+\Sigma_{es}\Sigma_{ss}^{-1}(s-\mu_s)=Ks+b,
\end{equation}
where
\begin{equation}
    K=\Sigma_{es}\Sigma_{ss}^{-1},
    \qquad
    b=\mu_e-K\mu_s.
\end{equation}
This derivation explains why an affine fast path is a plausible low-latency
approximation.  It also marks the boundary of the approximation.  Near lattice
decision boundaries, the posterior can be multimodal; maximum-likelihood,
closest-lattice-point, or wrapped-Gaussian decoders can outperform a single
global affine rule \cite{fukui2018,roy2025}.

\subsection{Logical failure criterion}

After correction, the residual displacement is wrapped into the GKP
fundamental cell.  A logical error occurs when the residual crosses a logical
decision boundary, for example
\begin{equation}
    |r_{q,t}|>\lambda/2 \quad \Rightarrow \quad X_L\ {\rm error},
\end{equation}
and similarly for the \(p\)-quadrature.  Closed-loop logical error probability
is therefore the central metric; parameter-regression error is useful only when
it correlates with the downstream correction outcome.

\section{Noise and Drift Model}

\subsection{Effective noise state}

The slow loop does not attempt to infer every microscopic hardware parameter.
It uses a low-dimensional effective state,
\begin{equation}
    \theta_t^{\rm noise}
    =
    (\sigma_t,\mu_{q,t},\mu_{p,t},\vartheta_t),
\end{equation}
where \(\sigma_t\) controls the displacement scale, \(\mu_{q,t}\) and
\(\mu_{p,t}\) represent mean bias, and \(\vartheta_t\) captures covariance-axis
rotation.  This compact parameterization can be mapped to affine runtime
parameters and estimated from finite histogram windows.

\subsection{Noise sources represented by the model}

The model is an effective displacement-noise model with additional measurement
uncertainty.  It absorbs several physical effects:
\begin{enumerate}
    \item finite-energy GKP peak width and envelope effects;
    \item Gaussian random displacement noise in \(q\) and \(p\);
    \item biased displacement means caused by calibration offsets;
    \item anisotropic or rotated covariance caused by asymmetric noise;
    \item syndrome measurement noise and noisy auxiliary-state contributions.
\end{enumerate}
This is narrower than a full circuit-level GKP noise model, but it is aligned
with the affine fast-path contract.

\subsection{Drift scenarios}

The benchmark suite uses four drift scenarios:
\begin{enumerate}
    \item \texttt{static\_bias\_theta}: a biased and rotated operating point;
    \item \texttt{linear\_ramp}: slowly changing effective noise parameters;
    \item \texttt{step\_sigma\_theta}: a sudden variance/orientation change;
    \item \texttt{periodic\_drift}: periodic nonstationarity.
\end{enumerate}
These scenarios separate different adaptation demands: steady calibration,
tracking, shock response, and periodic following.

\section{Model Architecture}

\subsection{Fast loop: affine fixed-point decoder}

The fast loop receives the current measured syndrome \(s_t\), reads the active
parameter bank, and computes
\begin{equation}
    s_t^{\rm clip}={\rm clip}(s_t,-s_{\max},s_{\max}),
\end{equation}
\begin{equation}
    \Delta_t^{\rm raw}=K_t s_t^{\rm clip}+b_t,
\end{equation}
\begin{equation}
    \Delta_t=Q\!\left({\rm clip}(\Delta_t^{\rm raw},
    -\Delta_{\max},\Delta_{\max})\right),
\end{equation}
where \(Q(\cdot)\) denotes fixed-point quantization.  The implementation uses
an FPGA-oriented Q4.20 representation for syndrome values and runtime
parameters.  The fast loop also records diagnostic counters, including
histogram-input saturation, correction saturation, overflow, aggressive
parameter events, commit counts, and fallback-related signals.

\subsection{Parameter mapping}

Given an estimated noise state, the parameter mapper constructs an error
covariance
\begin{equation}
    C =
    R(\vartheta)
    \begin{bmatrix}
        \sigma_q^2 & 0\\
        0 & \sigma_p^2
    \end{bmatrix}
    R(\vartheta)^\top ,
\end{equation}
and a measurement covariance
\begin{equation}
    R_{\rm meas}=(\sigma_{\rm meas}^2+\Delta_{\rm eff}^2)I.
\end{equation}
The raw gain is
\begin{equation}
    K_{\rm raw}=C(C+R_{\rm meas})^{-1},
\end{equation}
followed by gain clipping or scaling.  The bias target is
\begin{equation}
    b_{\rm target}=\alpha(I-K_{\rm target})\mu,
    \qquad
    \mu=\begin{bmatrix}\mu_q\\\mu_p\end{bmatrix}.
\end{equation}
Finally, exponential smoothing produces staged parameters:
\begin{equation}
    K_t=(1-\beta)K_{t-1}+\beta K_{\rm target},
    \qquad
    b_t=(1-\beta)b_{t-1}+\beta b_{\rm target}.
\end{equation}

\subsection{Teacher estimators}

The teacher family provides stable estimates of \(\theta_t^{\rm noise}\) from
recent syndrome history.  A simple teacher computes moments over a histogram
window.  Stronger teachers such as EKF, UKF, RLS, or particle-filter variants
add state-space assumptions over time.  Abstractly,
\begin{equation}
    \hat{\theta}_t^{\rm teacher}={\rm Teacher}(H_{1:t}),
\end{equation}
and
\begin{equation}
    (K_t^{\rm teacher},b_t^{\rm teacher})
    =
    {\rm ParamMapper}(\hat{\theta}_t^{\rm teacher}).
\end{equation}
The teacher keeps the learned branch near an interpretable calibration surface
and creates meaningful ablation baselines.

\subsection{CNN residual branch}

The CNN consumes a short context of normalized syndrome histograms and
histogram deltas:
\begin{equation}
    X_t^{\rm hist}=[H_{t-c+1},\ldots,H_t,\Delta H_{t-c+2},\ldots,\Delta H_t].
\end{equation}
Teacher-side features may include
\begin{equation}
    z_t^{\rm teacher}
    =
    (\hat{\theta}_t^{\rm teacher},
     K_t^{\rm teacher},
     b_t^{\rm teacher},
     \Delta b_t^{\rm teacher},\ldots).
\end{equation}
The gated branch narrows this to a small set of teacher-\(b\) and
teacher-\(\Delta b\) scalars.  The learned branch predicts
\begin{equation}
    \widehat{\delta b}_t=f_\phi(X_t^{\rm hist},z_t^{\rm teacher}),
\end{equation}
then clips the result:
\begin{equation}
    \delta b_t={\rm clip}(s_b\widehat{\delta b}_t,-b_{\max},b_{\max}).
\end{equation}
The committed parameter is
\begin{equation}
    K_t=K_t^{\rm teacher},
    \qquad
    b_t={\rm EMA}(b_t^{\rm teacher}+\delta b_t).
\end{equation}

\subsection{Statistical calibration branch}

The same affine runtime contract also supports a non-neural statistical
calibration branch.  This branch estimates a bounded bias correction directly
from recent syndrome statistics and emits runtime parameters through the same
clipping, smoothing, and stage-and-commit path.  Its role is both practical and
scientific: it tests whether the benefit comes from a general histogram-driven
calibration principle rather than from the CNN architecture alone.

\subsection{Stage-and-commit runtime contract}

The slow loop does not directly mutate active fast-loop parameters.  It stages
a candidate \((K,b)\) into an inactive bank and commits it at a safe epoch
boundary:
\begin{equation}
(K_t,b_t)=
\begin{cases}
(K^A,b^A), & t<t_{\rm commit},\\
(K^B,b^B), & t\ge t_{\rm commit}.
\end{cases}
\end{equation}
This contract allows stale-parameter effects, update latency, commit success,
rollback/fallback behavior, and fixed-point stability to be measured
separately from logical performance.

\section{Relationship to Existing Work}

\subsection{GKP and bosonic analog soft information}

GKP analog syndrome information has already been shown to be valuable.  Work
on analog QEC, surface-GKP codes, and concatenated bosonic-QLDPC decoding uses
continuous measurement information to set decoder weights or soft messages
\cite{fukui2018,noh2020,noh2022,raveendran2022,berent2024,borah2025}.  The
distinction here is that the analog statistics are compressed into
physical-layer affine GKP fast-path parameters, rather than into outer-code
matching weights or LDPC messages.

\subsection{Adaptive priors and syndrome-statistics estimation}

Adaptive-weight and syndrome-statistics methods estimate decoder weights or
noise parameters from measured data \cite{spitz2018,wagner2021,dgr2023}.
Calibrated decoders and decoder-prior optimization show that prior mismatch can
materially affect logical performance \cite{chen2022,sivak2024}.  This work
borrows that lesson, but maps the adaptation target to the GKP affine
parameters \((K,b)\) and constrains the update through a stage-and-commit
runtime contract.

\subsection{Learned low-latency QEC modules}

AI pre-decoders, high-accuracy neural decoders, and calibration-conditioned
FiLM decoders demonstrate a systems-oriented pattern: a learned module is most
useful when its input/output contract, latency role, and integration with a
classical decoder are explicit \cite{chamberland2026,stein2026,
google2023decoder}.  The learned module here is not the per-shot decoder.  It
is a slow calibration component that updates low-dimensional affine
parameters; the per-shot path remains fixed-point affine logic.

\subsection{Real-time and FPGA QEC decoders}

Real-time QEC decoders and FPGA-integrated systems have reported low-latency
hardware decoding in stabilizer-code settings \cite{caune2024,yang2026,
ziad2024,maurer2025}.  Real-time bosonic QEC has also been demonstrated in
experimental feedback systems \cite{sivak2023,lachance2024}.  These works set
the standard for full deployment claims: closed-loop timing, worst-case
latency, resource use, and hardware integration.  The contribution here is a
hardware-compatible GKP affine calibration contract and a software-HIL
evaluation of its drift-adaptive behavior.

\section{Experimental Setup}

\subsection{Software-HIL protocol}

The numerical experiments use a software-HIL protocol that preserves the fast
loop, slow loop, parameter-bank, and staged-commit interfaces while executing
the underlying experiment in software.  This protocol is sufficient for
comparing adaptive calibration laws under identical drift trajectories and
paired seeds.  It does not measure board-level latency or hardware resource
use.

\subsection{Scenarios and modes}

The four drift scenarios are:
\begin{center}
\small
\begin{tabular}{ll}
\toprule
Scenario key & Interpretation \\
\midrule
\texttt{static\_bias\_theta} & biased and rotated operating point \\
\texttt{linear\_ramp} & slowly changing effective noise parameters \\
\texttt{step\_sigma\_theta} & sudden variance/orientation change \\
\texttt{periodic\_drift} & periodic nonstationarity \\
\bottomrule
\end{tabular}
\end{center}

The main comparison includes filtering baselines, residual affine baselines,
the learned hybrid residual branch, and statistical calibration variants.  The
reported metric is \(\texttt{final\_ler\_mean}\), where lower values indicate
fewer logical failures under the fixed software-HIL protocol.

\section{Numerical Results and Benchmark Plan}

\subsection{Four-scenario affine benchmark}

Table~\ref{tab:five-mode-benchmark} compares five affine decoding modes under
the four drift scenarios.  The learned hybrid residual branch is the best mode
among these five in all four scenarios, with UKF as the strongest filtering
baseline.

\begin{table}[ht]
\centering
\small
\begin{tabular}{lrrrrr}
\toprule
Scenario & EKF & UKF & Const.-\(\mu\) & RLS-\(b\) & Hybrid-\(b\) \\
\midrule
\texttt{static\_bias\_theta} & 0.838110 & 0.825370 & 0.836658 & 0.837577 & \textbf{0.810902} \\
\texttt{linear\_ramp} & 0.819200 & 0.811201 & 0.816911 & 0.819373 & \textbf{0.787755} \\
\texttt{step\_sigma\_theta} & 0.822365 & 0.811548 & 0.819784 & 0.821493 & \textbf{0.788800} \\
\texttt{periodic\_drift} & 0.832192 & 0.821558 & 0.829670 & 0.832334 & \textbf{0.806392} \\
\bottomrule
\end{tabular}
\caption{Four-scenario affine benchmark.  Entries are
\(\texttt{final\_ler\_mean}\); lower is better.}
\label{tab:five-mode-benchmark}
\end{table}

\subsection{Feature and teacher ablations}

Table~\ref{tab:feature-ablation} isolates the learned residual branch.  Removing
histogram deltas degrades performance relative to the full hybrid model,
indicating that temporal changes in the syndrome histogram are useful for
tracking drift.  Removing teacher prediction also degrades the full hybrid
model, but the no-teacher-params variant performs best in this ablation set.
This result argues against a simple story in which every teacher-derived
channel is uniformly beneficial.

\begin{table}[ht]
\centering
\small
\begin{tabular}{lrrr}
\toprule
Mode & Avg. LER & \(\Delta\) vs UKF & \(\Delta\) vs Hybrid Full \\
\midrule
\texttt{ukf} & 0.817382 & 0.000000 & +0.018837 \\
\texttt{hybrid\_full} & 0.798545 & -0.018837 & 0.000000 \\
\texttt{hybrid\_no\_hist\_deltas} & 0.826723 & +0.009341 & +0.028178 \\
\texttt{hybrid\_no\_teacher\_prediction} & 0.807251 & -0.010131 & +0.008706 \\
\texttt{hybrid\_no\_teacher\_params} & \textbf{0.749621} & -0.067761 & -0.048924 \\
\texttt{hybrid\_no\_teacher\_deltas} & 0.800329 & -0.017053 & +0.001784 \\
\bottomrule
\end{tabular}
\caption{Feature and teacher ablation summary.  Negative
\(\Delta\) means improvement over the reference.}
\label{tab:feature-ablation}
\end{table}

\subsection{Statistical calibration as a strong comparator}

Table~\ref{tab:statcalib-extension} adds a statistical calibration lane to the
same four-scenario suite.  The statistical calibration branch beats both the
hybrid residual branch and UKF in every scenario.  This is an important result
for positioning the work: the central object is not only a CNN residual model,
but a broader adaptive affine calibration contract in which non-neural
statistical estimators can also be strong.

\begin{table}[ht]
\centering
\small
\begin{tabular}{lrrr}
\toprule
Scenario & UKF & Hybrid-\(b\) & StatCalib \\
\midrule
\texttt{static\_bias\_theta} & 0.825370 & 0.810902 & \textbf{0.431708} \\
\texttt{linear\_ramp} & 0.811201 & 0.787755 & \textbf{0.467083} \\
\texttt{step\_sigma\_theta} & 0.811548 & 0.788800 & \textbf{0.460016} \\
\texttt{periodic\_drift} & 0.821558 & 0.806392 & \textbf{0.438751} \\
\bottomrule
\end{tabular}
\caption{Statistical calibration comparison under the same drift
suite.  Entries are \(\texttt{final\_ler\_mean}\); lower is better.}
\label{tab:statcalib-extension}
\end{table}

\subsection{Local sensitivity of statistical calibration}

The statistical calibration branch was evaluated on a five-point local grid
around the default parameter setting: default scale/clip/threshold, low
residual scale, high residual scale, low residual clip, and high signal
threshold.  Table~\ref{tab:statcalib-sensitivity} shows that the advantage is
not tied to a single parameter point.  The high-threshold variant has the best
aggregate mean LER by a very small margin, whereas the default variant wins
three of four scenarios and has the best mean within-statcalib rank.

\begin{table}[ht]
\centering
\small
\begin{tabular}{lrrr}
\toprule
Variant & Mean LER & Mean Rank & Scenario Wins \\
\midrule
\texttt{statcalib\_high\_threshold} & \textbf{0.449241} & 1.75 & 1 \\
\texttt{statcalib\_default} & 0.449254 & \textbf{1.25} & \textbf{3} \\
\texttt{statcalib\_high\_scale} & 0.456477 & 3.00 & 0 \\
\texttt{statcalib\_low\_clip} & 0.481484 & 4.50 & 0 \\
\texttt{statcalib\_low\_scale} & 0.485129 & 4.50 & 0 \\
\bottomrule
\end{tabular}
\caption{Local sensitivity of statistical calibration.  Mean rank is
computed within the five statistical calibration variants for each scenario.}
\label{tab:statcalib-sensitivity}
\end{table}

\subsection{Per-scenario best statistical calibration variant}

Table~\ref{tab:statcalib-scenario-best} reports the best statistical
calibration variant in each scenario.  Every scenario-best variant clears both
the UKF and hybrid residual baselines by a large margin.  The
\texttt{static\_bias\_theta} winner is the only row with a mixed status flag,
so that scenario is best treated as a useful but locally caveated win.

\begin{table}[ht]
\centering
\small
\begin{tabular}{lrrr}
\toprule
Scenario & Best variant & Best LER & Gap vs Hybrid-\(b\) \\
\midrule
\texttt{static\_bias\_theta} & \texttt{high\_threshold} & 0.431684 & 0.379889 \\
\texttt{linear\_ramp} & \texttt{default} & 0.467094 & 0.321178 \\
\texttt{step\_sigma\_theta} & \texttt{default} & 0.458834 & 0.328883 \\
\texttt{periodic\_drift} & \texttt{default} & 0.439352 & 0.367232 \\
\bottomrule
\end{tabular}
\caption{Best statistical calibration variant per scenario.  Gap is
the hybrid residual LER minus the best statistical calibration LER.}
\label{tab:statcalib-scenario-best}
\end{table}

\subsection{Mechanism probe for residual-b behavior}

A multi-seed mechanism figure is available for the learned residual branch.  It
compares seed-wise hybrid-model gaps and a lower-clip intervention.  The
descriptive reading is that residual-b high-amplitude behavior appears in more
than one seed, but lowering the residual clip is mixed and often harmful.  This
supports a cautious interpretation: residual amplitude participates in some
failure modes, but it is not a monotonic causal explanation by itself.

\subsection{Unseen drift generalization}

The current four scenarios cover steady bias, ramp tracking, step response, and
periodic drift.  A stronger benchmark requires unseen drift families such as
random-walk drift, \(1/f\)-like drift, burst measurement noise, coupled
bias-variance drift, and faster-than-window drift.  This experiment is needed
to distinguish genuine drift adaptation from overfitting to a small set of
hand-designed trajectories.

\subsection{Oracle and wrapped-Gaussian baselines}

The affine methods need comparison against stronger theoretical baselines:
nearest-lattice hard decoding, static affine decoding, oracle affine decoding
with true noise state, and wrapped-Gaussian or maximum-likelihood decoding in
small settings.  These baselines quantify how much performance is lost because
the fast path is affine, and how much is lost because the slow loop estimates
the noise state imperfectly.

\subsection{Runtime, quantization, and fixed-point degradation}

The deployment claim depends on more than logical performance.  The benchmark
requires fixed-point versus floating-point comparisons, saturation and
overflow rates, commit latency, stale-parameter penalty, fallback frequency,
and per-update slow-loop cost.  These measurements connect the mathematical
affine rule to the real-time decoder contract.

\subsection{Embedded runtime and board-level validation}

Embedded inference and board-level validation are separate from software-HIL
simulation.  The relevant measurements are exported-model equivalence, embedded
runtime latency, host-to-device update cost, hardware resource use, and
closed-loop timing.  These experiments are needed before replacing
``hardware-compatible'' with a stronger hardware-demonstrated claim.

\section{Discussion}

\section{Conclusion}

This work frames drift-adaptive GKP correction as a two-timescale affine
calibration problem.  The fast loop remains a bounded fixed-point rule,
\(\Delta_t=K_t s_t+b_t\), while a slower loop updates the affine parameters
from recent syndrome statistics.  This formulation creates a practical middle
ground between fixed decoders, which cannot track hardware drift, and full
neural decoders, which can be difficult to place on a deterministic per-shot
path.

The numerical results support two main conclusions.  First, adaptive affine
decoding with a teacher-anchored residual branch improves over filtering
baselines in the four-scenario drift suite, and ablations show that temporal
histogram changes carry useful information.  Second, statistical calibration is
a strong comparator under the same runtime contract; its advantage persists
under a local sensitivity grid.  Together, these results suggest that the most
robust scientific story is not simply that a CNN is better than classical
filters, but that low-dimensional histogram-driven calibration of an affine GKP
fast path is a promising structure for drift-adaptive decoding.

The remaining technical gap is to connect this promising software-HIL evidence
to broader drift families, stronger oracle and wrapped-Gaussian baselines, and
measured runtime constraints.  Those additions would test whether the affine
calibration contract remains useful when the benchmark moves from a controlled
drift suite toward a deployment-grade decoder evaluation.

\begin{thebibliography}{99}

\bibitem{gkp2001}
D. Gottesman, A. Kitaev, and J. Preskill,
``Encoding a qubit in an oscillator,''
\emph{Physical Review A}, 2001.

\bibitem{grimsmo2021}
A. L. Grimsmo and S. Puri,
``Quantum Error Correction with the Gottesman--Kitaev--Preskill Code,''
\emph{PRX Quantum} 2, 020101, 2021.

\bibitem{fukui2018}
K. Fukui, A. Tomita, A. Okamoto, and K. Fujii,
``High-threshold fault-tolerant quantum computation with analog quantum error
correction,''
\emph{Physical Review X}, 2018.

\bibitem{noh2020}
K. Noh and C. Chamberland,
``Fault-tolerant bosonic quantum error correction with the surface-GKP code,''
\emph{Physical Review A}, 2020.

\bibitem{noh2022}
K. Noh, C. Chamberland, and F. G. S. L. Brandao,
``Low-overhead fault-tolerant quantum error correction with the surface-GKP
code,''
\emph{PRX Quantum}, 2022.

\bibitem{raveendran2022}
N. Raveendran, N. Rengaswamy, F. Rozpedek, A. Raina, L. Jiang, and B. Vasic,
``Finite-rate QLDPC-GKP coding scheme that surpasses the CSS Hamming bound,''
\emph{Quantum}, 2022.

\bibitem{berent2024}
L. Berent et al.,
``Analog information decoding of bosonic quantum LDPC codes,''
\emph{PRX Quantum}, 2024.

\bibitem{borah2025}
S. K. Borah, A. K. Pradhan, N. Raveendran, M. Pacenti, and B. Vasic,
``Fault tolerant decoding of QLDPC-GKP codes with circuit level soft
information,''
preprint, 2025.

\bibitem{roy2025}
M.-A. Roy, T. Pousset, and B. Royer,
``Decoding multimode Gottesman--Kitaev--Preskill codes with noisy auxiliary
states,''
preprint, 2025.

\bibitem{spitz2018}
S. Spitz, B. Tarasinski, C. W. J. Beenakker, and T. E. O'Brien,
``Adaptive weight estimator for quantum error correction,''
\emph{Advanced Quantum Technologies}, 2018.

\bibitem{wagner2021}
T. Wagner, H. Kampermann, and D. Bru{\ss},
``Optimal noise estimation from syndrome statistics of quantum codes,''
\emph{Physical Review Research}, 2021.

\bibitem{dgr2023}
H. Chen et al.,
``DGR: Tackling drifted and correlated noise in quantum error correction via
decoding graph re-weighting,''
preprint, 2023.

\bibitem{chen2022}
E. H. Chen et al.,
``Calibrated decoders for experimental quantum error correction,''
\emph{Physical Review Letters}, 2022.

\bibitem{sivak2024}
V. Sivak, M. Newman, and P. Klimov,
``Optimization of decoder priors for accurate quantum error correction,''
\emph{Physical Review Letters} 133, 150603, 2024.

\bibitem{chamberland2026}
C. Chamberland, J. Olle, M. Li, S. Thornton, and I. Baratta,
``Fast and accurate AI-based pre-decoders for surface codes,''
arXiv:2604.12841, 2026.

\bibitem{stein2026}
S. Stein, S. Kan, C. Liu, A. Harkness, S. Garner, Z. Du, Y. Ding, Y. Mao, and
A. Li,
``Calibration-conditioned FiLM decoders for low-latency decoding of quantum
error correction evaluated on IBM repetition-code experiments,''
preprint, 2026.

\bibitem{google2023decoder}
Google Quantum AI and collaborators,
``Learning high-accuracy error decoding for quantum processors,''
preprint, 2023.

\bibitem{caune2024}
L. Caune et al.,
``Demonstrating real-time and low-latency quantum error correction with
superconducting qubits,''
preprint, 2024.

\bibitem{yang2026}
X. Yang et al.,
``Real-time surface-code error correction using an FPGA-based neural-network
decoder,''
preprint, 2026.

\bibitem{ziad2024}
M. Ziad et al.,
``Local clustering decoder: a fast and adaptive hardware decoder for the
surface code,''
preprint, 2024.

\bibitem{maurer2025}
L. Maurer et al.,
``Real-time decoding of the gross code memory with FPGAs,''
preprint, 2025.

\bibitem{sivak2023}
V. Sivak et al.,
``Real-time quantum error correction beyond break-even,''
\emph{Nature}, 2023.

\bibitem{lachance2024}
D. Lachance-Quirion et al.,
``Autonomous quantum error correction of Gottesman--Kitaev--Preskill states,''
\emph{Physical Review Letters}, 2024.

\end{thebibliography}

\end{document}
